% ======================================================================
%  KAVRAM: Workflow  (Tier 1)  — kompakt
% ======================================================================

\concept{Orchestration · Orkestrasyon}{Workflow}{Adımları modelin değil senin tanımladığın, kontrollü orkestrasyon}

\stage{1}{Problem}
Çok adımlı ama \emph{adımları belli} bir AI görevin var --- yine de onu otonom
bir agent olarak kurdun. Şimdi maliyet öngörülemez, hata ayıklama zor, bazen
döngüde sıkışıyor. Oysa ihtiyacın esneklik değil, \textbf{güvenilirlik ve
kontroldü}.

\keyline{Adımları zaten biliyorsan, onlara modelin karar vermesine gerek yok. Karar verdiğin her şey kontrol kaybıdır.}

\stage{2}{Neden Var}
Adımlar belliyse orkestrasyonu sen sabitlersin: sıralı adımlar, açık dallanma,
ve model yalnızca belirli noktalarda belirli alt-görevler için çağrılır. Sonuç:
öngörülebilirlik, ucuzluk, ve adım adım test edilebilirlik.

\stage{3}{Tarihçe}
\begin{timeline}
\yr{2022} tek prompt \tlsep \yr{2023} prompt chaining: bilinen adımları zincirle \tlsep \yr{2023} orkestre boru hatları \tlsep \yr{2024–26} workflow vs agent ayrımı netleşir: kontrol vs otonomi
\end{timeline}

\stage{4}{Tanım}
\begin{problembox}
\textbf{Workflow:} Adımları ve kontrol akışını \emph{senin} tanımladığın sistem;
model belirli noktalarda belirli alt-görevler için çağrılır. Deterministik
orkestrasyon + olasılıksal adımlar. Agent'ın karşıtı: kararı modele değil,
koda verirsin.
\end{problembox}

\stage{5}{Zihinsel Model}
\begin{mentalbox}
Workflow bir \textbf{montaj hattıdır}: istasyonlar sabit. Model her istasyonda
tek bir işi yapan \emph{yetenekli işçidir} --- ama ürünün bir sonraki nereye
gideceğine \emph{bant} (senin kodun) karar verir, işçi değil. Öngörülebilir ve
denetlenebilir; karşılığında giderken yeni bir yol icat edemez.
\end{mentalbox}

\stage{6}{Mimari Diyagram}
Sabit adımlar + açık dallanma; kontrol akışı kodda, model düğümlerde.
\vspace{3pt}
\begin{center}
\begin{tikzpicture}[font=\sffamily\scriptsize, node distance=8mm,
   box/.style={draw=accent, line width=0.8pt, fill=wash, rounded corners=2pt,
               minimum height=9mm, minimum width=18mm, align=center},
   d/.style={draw=accent, line width=0.8pt, fill=wash, diamond, aspect=1.8,
             align=center, inner sep=0pt, minimum width=20mm},
   io/.style={draw=softgray, fill=white, rounded corners=8pt, align=center,
              minimum height=9mm, minimum width=14mm, inner sep=2pt},
   flow/.style={-{Stealth[length=2mm]}, accent, line width=0.8pt}]
  \node[io] (in) {Girdi};
  \node[box, right=8mm of in] (s1) {Çıkar\\\textcolor{softgray}{LLM}};
  \node[d, right=9mm of s1] (s2) {Sınıflandır};
  \node[box, above right=4mm and 9mm of s2] (s3a) {Özetle\\\textcolor{softgray}{LLM}};
  \node[box, below right=4mm and 9mm of s2] (s3b) {Yönlendir};
  \node[io, right=24mm of s2, fill=good!10, draw=good] (out) {Çıktı};
  \draw[flow] (in) -- (s1);
  \draw[flow] (s1) -- (s2);
  \draw[flow] (s2) -- node[above left=-1mm, font=\sffamily\tiny, text=accentdk]{A} (s3a);
  \draw[flow] (s2) -- node[below left=-1mm, font=\sffamily\tiny, text=accentdk]{B} (s3b);
  \draw[flow] (s3a) -| (out);
  \draw[flow] (s3b) -| (out);
\end{tikzpicture}
\end{center}
\keyline{Her ok senin kodundur, modelin kararı değil. Bu yüzden her adım ayrı test edilebilir, maliyet öngörülebilir.}

\stage{7}{Karar Çerçevesi}
Agent kararının diğer yüzü. Hangi işaretler ne der?\par
\vspace{2pt}
\noindent\begin{tabularx}{\linewidth}{@{}L L@{}}
\arrayrulecolor{hairline}\toprule
{\sffamily\bfseries\small Workflow seç} & {\sffamily\bfseries\small Agent seç} \\
\midrule
Adımlar önceden belli & Adımlar çalışırken belli olur \\
\addlinespace[2pt]
Öngörülebilirlik / denetim kritik & Açık uçlu, uyarlanır görev \\
\addlinespace[2pt]
Maliyet sabit olmalı & Esneklik değeri riskten yüksek \\
\addlinespace[2pt]
Çoğu üretim akışı & Gerçek belirsizlik (azınlık) \\
\bottomrule
\end{tabularx}
\keyline{Varsayılan Workflow. Çoğu ``agent'' aslında birkaç dallanmalı bir workflow'dur.}

\stage{8}{Üretim Senaryosu}
\begin{prodbox}[Üretim Senaryosu · Belge İşleme Hattı]
\textbf{Bağlam.} Gelen belgeler işlenecek: bilgi çıkar, türünü sınıflandır,
özetle, doğru ekibe yönlendir. İlk sürüm tek bir agent'tı --- güvenilmez ve
hata ayıklaması zordu.

\smallskip
\textbf{Çözüm.} Sabit bir workflow: \emph{çıkar → sınıflandır → (türüne göre)
özetle/yönlendir}. Her adım modelle ama ayrı, test edilebilir bir düğüm.

\smallskip
\textbf{Sonuç.} Öngörülebilir maliyet, adım adım eval, kolay hata ayıklama.
Esneklik kaybı yok --- çünkü adımlar zaten baştan belliydi.
\end{prodbox}

\stage{9}{Anti-pattern}
\begin{antibox}[Anti-pattern · ``Mega-Prompt'' (Adımlara Bölmemek)]
Çok adımlı bir işi tek bir devasa prompt'a sıkıştırmak --- çıkar, sınıflandır,
özetle, karar ver; hepsi tek çağrıda. Sonuç: güvenilmez, test edilemez, bir
adım bozulduğunda nerede bozulduğu görünmez.

\smallskip
\textbf{Belirti:} ``Dev prompt bazen çalışıyor, neden bozulduğunu bilmiyoruz.''\\
\textbf{Kök neden:} İş, ayrık ve test edilebilir adımlara bölünmemiş.\\
\textbf{Düzeltme:} İşi bir workflow'a böl; her adımı ayrı prompt + ayrı eval yap.
\end{antibox}

\stage{10}{Alternatifler \& Ödünleşimler}
\vspace{2pt}
\begin{tabularx}{\linewidth}{@{}>{\raggedright\arraybackslash}p{26mm} L L@{}}
\arrayrulecolor{hairline}\toprule
{\sffamily\bfseries\small Yaklaşım} & {\sffamily\bfseries\small Güçlü} & {\sffamily\bfseries\small Ödünleşim} \\
\midrule
\textbf{Workflow} & Öngörülebilir, ucuz, denetlenir & Sabit; açık uçluya uymaz \\
\addlinespace[2pt]
\textbf{Agent} & Açık uçlu görev & Öngörülemez; pahalı; zor debug \\
\addlinespace[2pt]
\textbf{Prompt chain} & Bilinen düz adımlar & Dallanma/karar zayıf \\
\addlinespace[2pt]
\textbf{Tek prompt} & En basit iş & Çok adımlıya yetmez \\
\bottomrule
\end{tabularx}

\stage{11}{İlişkili Kavramlar}
\term{Agent} (karşıtı) \tlsep \term{Prompt Chaining} (basit workflow) \tlsep
\term{Router} (dallanma noktası) \tlsep \term{Tool Calling} (adımlardaki eylem) \tlsep
\term{Eval} (adım adım test)

\stage{12}{Özet}
\begin{itemize}
  \item Workflow: adımları ve akışı sen tanımlarsın; model düğümlerde çalışır.
  \item Öngörülebilir, ucuz, adım adım test edilebilir --- agent'ın veremediği kontrol.
  \item Adımlar önceden belliyse workflow; ancak gerçek belirsizlikte agent.
  \item Çok adımlı işi tek mega-prompt'a sıkıştırma; düğümlere böl.
  \item Çoğu üretim akışı, birkaç dallanmalı bir workflow'dur.
\end{itemize}

\stage{13}{Mühendislik İlkesi}
\begin{principlebox}
\textbf{İlke \#12 — Adımları biliyorsan onları sen yaz; modele yalnızca dolduramadığın boşlukları bırak.}\\[4pt]
\normalsize Kontrol bir maliyet değil, bir armağandır. Modele yalnızca gerçekten
karar gerektiren yeri devret.
\end{principlebox}
